\documentclass[11pt,a4paper]{article}
\usepackage[margin=27mm]{geometry}
\usepackage{fontspec}
\setmainfont{Latin Modern Roman}
\usepackage{amsmath,amssymb,amsthm,mathtools}
\usepackage{microtype}
\usepackage{xcolor}
\definecolor{InternalLink}{HTML}{245B53}
\definecolor{CitationLink}{HTML}{7A4937}
\definecolor{ExternalLink}{HTML}{365A83}
\usepackage[colorlinks=true,linkcolor=InternalLink,citecolor=CitationLink,urlcolor=ExternalLink]{hyperref}
\usepackage[nameinlink,noabbrev]{cleveref}
\newtheorem{theorem}{Theorem}[section]
\newtheorem{proposition}[theorem]{Proposition}
\newtheorem{lemma}[theorem]{Lemma}
\newtheorem{corollary}[theorem]{Corollary}
\theoremstyle{definition}
\newtheorem{definition}[theorem]{Definition}
\newtheorem{problem}[theorem]{Problem}
\theoremstyle{remark}
\newtheorem{remark}[theorem]{Remark}
\newcommand{\Q}{\mathbf Q}
\newcommand{\Z}{\mathbf Z}
\newcommand{\CH}{\operatorname{CH}}
\newcommand{\Pic}{\operatorname{Pic}}
\newcommand{\rk}{\operatorname{rk}}
\newcommand{\ch}{\operatorname{ch}}
\newcommand{\Td}{\operatorname{Td}}
\newcommand{\cP}{\mathcal P}
\newcommand{\cO}{\mathcal O}
\newcommand{\cD}{\mathcal D}
\newcommand{\MathEqref}[1]{\eqref{#1}}
\newcommand{\MathSecref}[1]{\Cref{#1}}
\setlength{\emergencystretch}{3em}
\title{Normalized Determinants and Height Metrics}
\author{}
\date{17 November 2026}
\hypersetup{pdftitle={Normalized Determinants and Height Metrics},pdfauthor={}}
\begin{document}
\maketitle
\begin{abstract}
We compare normalized determinant-of-cohomology metrics on Jacobian
families. For a perfect complex on the product of a pointed curve and
a smooth projective variety, the derived projection formula, for product
Kähler metrics, induces an exact isometry after normalization at the origin. In virtual rank zero,
the resulting line is the Deligne pairing with the determinant of the
pushforward. Filtered Grothendieck--Riemann--Roch identifies the line
attached to a curve correspondence and a homologically trivial cycle.
At a finite place where the pointed source curve has good reduction,
a model-perfect complex of generic virtual rank zero and degree-zero
determinant yields the canonical rigidified Poincar\'e model metric.
\end{abstract}
\tableofcontents

\section{Introduction}
The comparison between a cycle-theoretic height and a determinant
construction has two parts. One must identify the underlying line, and
one must identify its metric with the correct normalization. The second
part contains a genuine constant ambiguity even when the two curvature
forms agree. A Jacobian family provides a useful setting in which that
ambiguity can be removed exactly.

Let \(C\) be a pointed smooth projective complex curve and
\(J=\Pic^0(C)\). The standard Poincar\'e metric on \(C\times J\)
has purely mixed curvature. In a determinant family obtained by
twisting a fixed complex with the Poincar\'e bundle, the family index
formula therefore depends only on the virtual rank of its pushforward
to \(C\). Dividing by the determinant fibre at the origin equips the
line with a unit rigidification. Equality of curvature forms then gives
equality of metrics: their logarithmic ratio is pluriharmonic on the
compact connected variety \(J\), and its value at the origin is zero.

This argument proves the projection-formula isometry in
\Cref{thm:quillen-isometry}. Its rank-zero specialization is
\Cref{thm:deligne-isometry}. Applied to a curve correspondence, the two
theorems give \Cref{cor:cycle-determinant}; the relevant rank and
determinant are computed by the leading terms of
Grothendieck--Riemann--Roch. The good-reduction analogue is
\Cref{prop:good-reduction}. There the quotient of two model lines has a
vertical divisor, and the unit rigidification forces its coefficient
to vanish.

The full arguments appear in the next section in this order. Throughout,
\(\CH^j(X)_{\Q}\) denotes the rational Chow group, and
\[
 V^j(X)=\ker\bigl(\CH^j(X)_{\Q}\longrightarrow
                    H^{2j}(X,\Q)\bigr)
\]
over \(\mathbf C\). The cochain determinant convention and the
Poincar\'e normalization are fixed explicitly before the comparison
theorems. The determinant--biextension identifications of
\cite{Gorchinskiy,Nakayama} explain how this comparison enters local
height formulas.


% Source: research-sources/height-r2/sections/determinants.tex
\section{Normalized determinants on Jacobian families}\label{sec:metrics}

\subsection{The Poincaré metric and normalization}\label{sec:poincare-metric}

Let \(C\) be a smooth projective connected complex curve with
basepoint \(o\), and let \(J=\Pic^0(C)\). Let \(\cP\) be the
Poincaré line bundle on \(C\times J\), rigidified on both axes.
Equip it with the standard smooth metric for which these
rigidifications have norm one. Its Chern form
\(\Omega=c_1(\cP)\) is purely mixed: it has real degree one in
\(C\) and degree one in \(J\).

One way to see this is to pull back the canonical Poincaré bundle
on \(J\times J^\vee\) by the Abel map and the principal
polarization, with the convention realizing \(\cP\). For a
translation-invariant theta curvature \(\theta\), the corresponding
Poincaré curvature is the mixed form obtained from
\(m^*\theta-p_1^*\theta-p_2^*\theta\), up to the chosen duality
sign. The pure terms cancel. Pullback by the Abel map preserves
the mixed-degree assertion.

For a fixed smooth projective complex variety \(Z\), use the
cochain determinant convention
\[
 \lambda_Z(K)=\det R\Gamma(Z,K)
       =\bigotimes_i\det H^i(Z,K)^{(-1)^i}.
\]
A perfect complex is given a bounded vector-bundle representative
with smooth hermitian metrics; its determinant metric is the
alternating tensor product of the Quillen metrics of the terms.
Choose fixed Kähler metrics on the varieties involved.

\subsection{The normalized projection-formula isometry}
\label{sec:quillen-proof}

Let \(X\) be a smooth projective connected complex variety,
\(M=C\times X\), and \(q:M\to C\). For a hermitian perfect
complex \(\mathcal F\) on \(M\), put \(K=Rq_*\mathcal F\)
and choose any hermitian vector-bundle representative of \(K\).
Define line bundles on \(J\) by
\begin{align*}
 \cD_M(\mathcal F)
  &=\det Rp_{J*}(\cP_M\otimes\mathcal F_J)
     \otimes\bigl(\lambda_M(\mathcal F)\otimes\cO_J\bigr)^{-1},\\
 \cD_C(K)
  &=\det Rp_{J*}(\cP\otimes K_J)
     \otimes\bigl(\lambda_C(K)\otimes\cO_J\bigr)^{-1}.
\end{align*}
The subscripts indicate pullbacks to the appropriate products with
\(J\). Each constant denominator has the Quillen norm of the same
complex and metrics that occur in the numerator at \(0\).
Thus each line has a canonical origin rigidification of norm one.

\begin{theorem}\label[theorem]{thm:quillen-isometry}
Use product Kähler metrics on \(C\times X\). The derived
projection-formula isomorphism gives an exact isometry of
rigidified line bundles
\[
 \cD_M(\mathcal F)\xrightarrow{\sim}\cD_C(Rq_*\mathcal F).
\]
Both Chern forms equal
\begin{equation}\label{eq:quillen-curvature}
 \frac{\rk K}{2}\int_C\Omega^2.
\end{equation}
Under the canonical rigidified identifications, these normalized
metrics are independent of the hermitian representatives and the
fixed product Kähler metrics used to construct them.
\end{theorem}

\begin{proof}
Derived projection formula and base change identify the numerator
complexes and their restrictions at \(0\). Applying determinants
to these two identifications gives the asserted algebraic
isomorphism, respecting the origin rigidifications.

The Bismut--Gillet--Soulé curvature formula
\cite[Theorem~0.1]{BGS}, in the cochain determinant convention,
gives the Chern form of the first line as the base-degree-two
component of
\[
 \int_{C\times X}\Td(T_C)\Td(T_X)\ch(\mathcal F)
                          (e^\Omega-1).
\]
Every factor except \(\Omega\) is pulled back from \(C\times X\).
Since \(\Omega\) is purely mixed, only \(\Omega^2/2\) can
contribute degree two on \(J\). It already has degree two on
\(C\), so the remaining factor contributes its degree-zero-in-\(C\)
part. Integration over \(X\) gives
\[
 \int_X\Td(T_X)\ch(\mathcal F|_{\{x\}\times X})
       =\chi(X,\mathcal F|_{\{x\}\times X})=\rk K
 \qquad(x\in C).
\]
The final equality follows by derived base change and the constant
Euler characteristic of the perfect pushforward. This proves
\MathEqref{eq:quillen-curvature} for \(\cD_M\). The curve
calculation gives the same form: all positive-degree forms from
\(K\) or \(T_C\) have too much \(C\)-degree to contribute.

The determinant lines with their Quillen metrics are smooth on
the whole \(J\), including cohomology-jump loci. Under the
algebraic isomorphism, the logarithm of the ratio of their metrics
is therefore a smooth pluriharmonic function on the compact
connected manifold \(J\). It is constant, and the norm-one
origin rigidifications make that constant zero. Changing the
hermitian representatives or the fixed product metrics preserves both the
curvature computation and this normalization, proving the last
assertion.
\end{proof}

\subsection{Virtual rank zero and Deligne pairings}
\label{sec:deligne-proof}

\begin{theorem}\label[theorem]{thm:deligne-isometry}
Let \(K\) be a perfect complex of virtual rank zero on \(C\),
and put \(L=\det K\). There is a unique rigidified isomorphism
\[
 \cD_C(K)\xrightarrow{\sim}
       \langle\cP,p_C^*L\rangle_{C\times J/J}.
\]
It is an isometry for the normalized Quillen metric and the
Deligne metric. Both metrics are smooth and flat. The right side
can be formed using any smooth hermitian metric on \(L\).
\end{theorem}

\begin{proof}
On a nonsingular curve, rank and determinant identify
\(K_0(C)\) with \(\Z\oplus\Pic(C)\). Thus
\([K]=[L]-[\cO_C]\). At the level of line-bundle isomorphism
classes, the normalized determinant becomes
\[
 \lambda_C(\cP L)\lambda_C(\cP)^{-1}
       \lambda_C(L)^{-1}\lambda_C(\cO_C),
\]
the determinant expression for the Deligne pairing. Here the first
two factors are relative determinants over \(J\). Normalize an
isomorphism at \(0\). It becomes unique because an invertible
holomorphic function on the connected projective \(J\) is
constant.

The Quillen curvature is zero by
\MathEqref{eq:quillen-curvature}. The Deligne curvature is
\(\int_C\Omega\,c_1(L)=0\), since the integrand has degree
three on \(C\). At \(0\), \(\cP\) is the trivially metrized
line with unit rigidification; the Deligne unit has norm one.
The two flat metrics therefore agree by the same compactness
argument as above. The determinant and Deligne curvature formulas
are also discussed in \cite[Section~2]{BiswasSchumacher}.
\end{proof}

\subsection{Curve correspondences in arbitrary codimension}
\label{sec:correspondence-metric}

Let \(\dim X=n\), take
\(\Gamma\in\CH^r(C\times X)_\Q\), and let
\(W\in V^{n+1-r}(X)\). Choose rational \(K_0\)-lifts
\(E_\Gamma,E_W\) whose Chern characters start with \(\Gamma,W\)
in their indicated codimensions. Virtual sums of structure sheaves
give such lifts on these smooth varieties. Set
\[
 \mathcal F=E_\Gamma\otimes^{\mathbf L}p_X^*E_W,
 \qquad K=Rq_*\mathcal F.
\]
The product Chern character starts in codimension \(n+1\).
Filtered Grothendieck--Riemann--Roch for \(q\) gives
\begin{equation}\label{eq:grr-leading}
 \rk K=0,\qquad
 c_1(\det K)=q_*(\Gamma\cdot p_X^*W)={}^t\Gamma_*W.
\end{equation}
Positive-degree Todd terms do not contribute to these two
components. The divisor class on the right has degree zero because
\(W\) is homologically trivial.

\begin{corollary}\label[corollary]{cor:cycle-determinant}
There is a rigidified isometry of rational metrized line bundles
on \(J\),
\[
 \cD_{C\times X}(E_\Gamma\otimes^{\mathbf L}p_X^*E_W)
   \xrightarrow{\sim}
 \langle\cP,\cO_C({}^t\Gamma_*W)\rangle.
\]
The metric is the normalized flat Poincaré metric attached to the
degree-zero divisor class \({}^t\Gamma_*W\).
\end{corollary}

\begin{proof}
Combine \Cref{thm:quillen-isometry,thm:deligne-isometry} with
\MathEqref{eq:grr-leading}. For rational classes, take a common
tensor power throughout and then the positive real root of the
norm. Changing the chosen cycle lifts by higher filtration changes
their product only in codimension at least \(n+2\), which is
zero on \(C\times X\); the rational Chern character then gives
the same normalized line.
\end{proof}

The relation to the classical archimedean height uses the
determinant--cycle--Hodge biextension comparison of
\cite[Propositions~4.3, 4.29, 4.31 and Corollary~4.32]{Gorchinskiy}.
Its disjoint-cycle trivialization gives the local height
\cite[Remark~4.7]{Gorchinskiy}. The relative comparison over a
smooth base is given in \cite[Theorem~10.9]{Nakayama}; that theorem
includes an inversion of the central \(\mathbf G_m\) under its
conventions. Fixing the cochain determinant convention above and
the corresponding Poincaré convention fixes this sign.

An incidence locus can be a zero or pole locus of the chosen
meromorphic frame without being a singularity of the line metric.
If \(s=fe\), where \(e\) is a nonvanishing holomorphic frame, then
\[
 -\log\|s\|=-\log|f|-\log\|e\|.
\]
The second term is smooth across incidence in the Jacobian family
just constructed. This global line comparison resolves the
normalization and continuity issue for that determinant model.

\subsection{The source-curve good-reduction comparison}
\label{sec:good-reduction-metric}

Let \(\mathcal O\) be a complete discrete valuation ring with
fraction field \(K_v\), equipped with a fixed non-Archimedean
absolute value inducing the given discrete valuation.
Let \(\mathcal C/\mathcal O\) be a smooth proper geometrically
connected curve with a section, let \(\mathcal J/\mathcal O\)
be its Jacobian, and let \(\mathcal P\) be the axis-rigidified
Poincaré bundle. For a perfect complex \(K\) on \(\mathcal C\),
form the normalized determinant line \(\cD_{\mathcal C}(K)\)
on \(\mathcal J\) by the same algebraic construction.

\begin{proposition}\label[proposition]{prop:good-reduction}
Suppose the generic virtual rank of \(K\) is zero and
\(\deg\det K_{K_v}=0\). The model metric of
\(\cD_{\mathcal C}(K)\) on the generic fibre is the canonical
rigidified Poincaré metric associated with \(\det K_{K_v}\).
\end{proposition}

\begin{proof}
The degree of \(\det K\) is zero on every fibre. After
rigidification along the section, it defines a point of the
relative \(\Pic^0\), and hence a rigidified Poincaré extension
on \(\mathcal J\). The rank-zero determinant identity on the
generic curve identifies its generic line with
\(\cD_{\mathcal C}(K)\), respecting the origin rigidification.

Regard this isomorphism as a rational section of the quotient of
the two model lines. Its divisor is vertical. The abelian scheme
\(\mathcal J\) is regular and its special fibre is geometrically
integral, so this divisor is \(a[\mathcal J_s]\) for an integer
\(a\). Pullback to the zero section has valuation \(a\), but
the two integral unit rigidifications make that valuation zero.
Thus \(a=0\). Normality extends the isomorphism and its inverse
over the whole model. The model metrics coincide. The Poincaré
multiplication laws extend over the abelian scheme and are
isometries, identifying this model metric with the canonical one.
\end{proof}

If an ambient model \(\mathcal X\) and a model-perfect complex
on \(\mathcal C\times\mathcal X\) have a perfect derived
pushforward to \(\mathcal C\), the determinant projection
formula reduces their model comparison to this proposition.
At a bad-reduction source curve the vertical correction involves
the geometry of a semistable or Néron model; this is a separate
comparison problem.



\section{The remaining metric comparisons}\label{sec:remaining-metrics}

The comparisons above provide the complex-place and good-reduction
inputs for a local height comparison. The next step is to identify the
vertical correction at a place of semistable reduction.

\begin{problem}[Semistable source curves]\label[problem]{prob:semistable-metric}
Fix a complete discretely valued field \(K_v\) with its original
absolute value, a pointed smooth projective geometrically connected
curve \(C/K_v\), and a semistable model \(\mathcal C\). Let \(K\)
be a perfect complex on the model with generic virtual rank zero and
\(\deg\det K_{K_v}=0\). Construct the admissible correction to its
normalized determinant model metric and identify the resulting
rigidified metric with the canonical Poincar\'e metric attached to
\(\det K_{K_v}\). For a determinant arising from a curve
correspondence, compare this metric with the actual-form Green
construction, retaining the finite local lengths, prescribed cycle
representatives and native finite-extension degree.
\end{problem}

In the smooth-model case, \Cref{prop:good-reduction} supplies this
comparison. At semistable reduction the correction must account for
the components and intersection graph. Passage to an actual-form
height also uses its trace, correspondence covariance and
normalization on the complete actual test space.

The corresponding family problem concerns both the boundary and the
growth of the resulting height. One can formulate a precise target
without asserting that every family has positive height growth.

\begin{problem}[Height domination in a family]\label[problem]{prob:family-height}
Let \(F\) be a number field, \(S/F\) a smooth quasi-projective
variety, and \(f:\mathcal X\to S\) a smooth projective family of
relative dimension \(2r-1\). Let \(Z\) be a relative codimension
\(r\) cycle whose fibres are homologically trivial, and assume the
fibrewise pairings \(h_s\) have been constructed with the fixed
absolute height convention, invariant under extension of number fields.
Determine geometric hypotheses on \((f,Z)\)
under which the function
\[
 q_Z(s)=(-1)^r h_s(Z_s,Z_s)
\]
is the height of a canonically metrized biextension line, has controlled
boundary asymptotics on a specified projective compactification \(\overline S\),
and satisfies
\[
 q_Z(s)\ge\varepsilon h_A(s)-C\qquad(s\in U(\overline F))
\]
for a nonempty open \(U\subset S\), an ample line bundle \(A\)
on \(\overline S\), and constants \(\varepsilon>0\), \(C\).
\end{problem}

Such an inequality would give finiteness of bounded-height points of
bounded degree on \(U\) by the classical Northcott theorem for
\(h_A\). It is a statement about the parameter space, distinct from
Northcott for all integral classes in one fixed Chow group.
The relative biextension comparison over a smooth base is
\cite[Theorem~10.9]{Nakayama}. Its singular-fibre extension in
Theorem~11.2 requires a modified cycle with a specified cohomological
property; Lemma~11.1 obtains that modification under the Hodge
conjecture. These conditions belong to the boundary comparison in
\Cref{prob:family-height}.

\begin{thebibliography}{9}
\bibitem{BGS}
J.-M. Bismut, H. Gillet and C. Soul\'e,
\emph{Analytic torsion and holomorphic determinant bundles I:
Bott--Chern forms and analytic torsion},
Communications in Mathematical Physics \textbf{115} (1988), 49--78.
\href{https://www.imo.universite-paris-saclay.fr/~jean-michel.bismut/Bismut/1988.pdf}{Author-hosted article}.

\bibitem{BiswasSchumacher}
I. Biswas and G. Schumacher,
\emph{Deligne pairing and Quillen metric},
\href{https://arxiv.org/abs/1501.02539}{arXiv:1501.02539}.

\bibitem{Gorchinskiy}
S. Gorchinskiy, \emph{Notes on the biextension of Chow groups},
\href{https://arxiv.org/abs/0802.1437v2}{arXiv:0802.1437v2}.

\bibitem{Nakayama}
Y. Nakayama, \emph{Asymptotics of local height pairing},
\href{https://arxiv.org/abs/2512.22551v1}{arXiv:2512.22551v1},
27 December 2025.
\end{thebibliography}
\end{document}
